\section{Introduction}
\label{sec:intro}

Diabetic retinopathy (DR) is a microvascular complication of diabetes mellitus and remains a leading cause of preventable vision impairment among working-age adults worldwide \citep{teo2021,saeedi2019}. A recent meta-analysis estimated that about one in five people living with diabetes (22.3\%) has some degree of retinopathy, and projected that the absolute burden will keep increasing through the middle of this century as the prevalence of diabetes rises \citep{teo2021}. Vision loss from DR is largely avoidable when disease is found early and treated on time, which is why periodic fundus photography is recommended for every person with diabetes. The bottleneck is not the camera but the reader: grading millions of photographs per year requires trained graders and ophthalmologists who are scarce in exactly those regions where diabetes is growing fastest. This has made automated DR screening one of the most extensively studied applications of deep learning in medical imaging, from the first large-scale validations \citep{gulshan2016,ting2017} to prospective trials of autonomous systems in primary care \citep{abramoff2018}.

Clinically, DR is graded from a small set of visible lesions. Microaneurysms (\ma), intraretinal haemorrhages (\he), hard exudates (\ex) and soft exudates or cotton-wool spots (\se) are the hallmarks of non-proliferative disease, and neovascularisation marks the proliferative stage. The International Clinical Diabetic Retinopathy (ICDR) severity scale \citep{wilkinson2003}, a simplification of the ETDRS classification \citep{etdrs1991}, defines five grades --- no DR, mild, moderate, severe non-proliferative DR and proliferative DR --- by stating which lesions must be present and, for the severe stage, how widely they must be distributed (the well-known ``4--2--1'' rule: haemorrhages in four quadrants, venous beading in two, or intraretinal microvascular abnormalities in one). Three machine-learning tasks therefore describe the same clinical reasoning at different granularity: \emph{grading} (image level, five ordered classes), \emph{lesion segmentation} (pixel level) and \emph{lesion detection} (instance level).

In the literature, however, these three tasks are almost always solved separately. Grading networks are trained on image-level labels and are asked to rediscover lesions implicitly, which makes them data-hungry and hard to explain. Segmentation networks are trained on a few hundred pixel-annotated images and are scored by overlap measures that say little about whether a clinician could use the output. Detection networks are borrowed from natural-image object detection and struggle with objects that occupy a handful of pixels. Multi-task models that share an encoder do exist \citep{foo2020multitask,playout2019,zhou2019collab}, as do lesion-aware grading models \citep{wang2017zoomin,sun2021lat,huang2021lesioncl}, but in most of them the link between lesion evidence and the final grade is learned implicitly through shared weights. Nothing prevents such a model from predicting \emph{severe} DR for an image in which it finds no haemorrhage at all, and the practitioner has no principled way to know when to distrust the output.

\paragraph{Gaps addressed in this work.} We identify four specific gaps. \textbf{(G1) Lesion--grade inconsistency.} Grade predictions are not required to agree with the lesion evidence produced by the same system, so internally contradictory outputs are possible and are rarely measured. \textbf{(G2) Ordinal structure and label noise.} DR grades are ordered, and inter-grader disagreement concentrates on adjacent grades, yet cross-entropy over nominal classes ignores both facts. \textbf{(G3) Small-lesion sensitivity.} Microaneurysms are a few pixels wide in a megapixel image; losses and architectures that are fine for organs are blind to them. \textbf{(G4) Absence of principled referral.} Most systems output a single hard label with a miscalibrated confidence; very few provide a finite-sample statement about how often the clinically true grade is contained in the output, which is what a referral workflow needs.

\paragraph{Our approach.} We propose \mname, a lesion-grounded grading framework built on top of a lesion-segmentation network. The segmentation network produces four lesion-probability maps. A differentiable, connected-component-free \emph{lesion-evidence bottleneck} converts those maps into soft per-zone lesion statistics (counts at several confidence levels, number of local peaks and maximum probability for each of four lesion types in each of six retinal zones). Each (lesion, zone) statistic becomes a token; a small transformer lets a grade query attend jointly to a global image embedding and to these 24 evidence tokens, so that the grade is computed \emph{from} the lesions the same network sees. A \emph{clinical-rule consistency loss} encodes ICDR-style necessary conditions (``grade $\ge$ severe requires haemorrhage in all four quadrants or heavy soft exudates'') as soft logic with a product t-norm, with thresholds calibrated on the training labels, and penalises predictions that exceed what the lesion evidence permits. The grade head is \emph{ordinal} with Gaussian-smoothed soft targets that reflect adjacent-grade disagreement, and the resulting class probabilities are temperature-scaled and wrapped by \emph{split-conformal prediction} to give referral sets with a finite-sample coverage guarantee. On the lesion side, a size-aware focal-Tversky objective with a full-resolution shallow branch is used to keep microaneurysms visible, and the very same probability maps are converted to detection boxes by connected-component proposals, so that segmentation and detection share one set of weights.

\paragraph{Experimental scope and honesty about compute.} All experiments reported here were produced on a single four-core CPU workstation without a GPU (Section~\ref{sec:setup}). This is a deliberate design constraint rather than an incidental one: it forces the study to be explicit about what is learned end-to-end (the lesion segmenter) and what is learned on top of frozen representations (the grading heads), and it makes all reported numbers reproducible without specialised hardware. It also means that absolute grading accuracy is below what an end-to-end fine-tuned, high-resolution, GPU-trained model could reach, and we do not claim state-of-the-art accuracy on the DDR benchmark. What we do claim, and test with controlled ablations, multiple seeds and confidence intervals, is that \emph{given the same frozen image embedding}, grounding the grade in lesion evidence and enforcing rule consistency yields measurable gains in agreement, internal consistency, calibration and referral quality over the standard alternatives.

\paragraph{Contributions.}
\begin{enumerate}
\item We formulate DR grading, lesion segmentation and lesion detection as a single coupled pipeline in which the grade is a function of differentiable, zone-wise lesion statistics (Section~\ref{sec:evidence}), and we make this coupling explicit with a token-based evidence head.
\item We introduce a soft-logic \emph{clinical-rule consistency} regulariser with learnable, label-calibrated thresholds and a matching \emph{rule-violation rate} metric that quantifies internal contradictions between lesion evidence and predicted grade (Section~\ref{sec:rule}).
\item We combine ordinal cumulative-link learning with Gaussian-smoothed soft targets and split-conformal prediction to obtain calibrated grade distributions and referral sets with finite-sample coverage (Sections~\ref{sec:ordinal} and~\ref{sec:conformal}).
\item We provide a size-aware lesion segmentation recipe and a segmentation-derived detection protocol with FROC analysis for lesions that are only a few pixels wide, evaluated on the full DDR lesion subset (Section~\ref{sec:seg}).
\item We report a comprehensive and fully reproducible empirical study on the DDR dataset \citep{li2019ddr} --- 13{,}673 fundus images, 757 of them with pixel masks and boxes --- including loss and preprocessing ablations, component ablations over five seeds, calibration and conformal analyses, and a zero-shot external test on the Gaussian-filtered APTOS-derived Kaggle set \citep{aptos2019,rath2020gaussian} (Section~\ref{sec:results}). All code and cached lesion maps are released.
\end{enumerate}

\paragraph{Organisation.} Section~\ref{sec:related} reviews related work and positions \mname\ against it. Section~\ref{sec:data} describes the datasets and preprocessing. Section~\ref{sec:method} details the method, with mathematical formulation and algorithms. Section~\ref{sec:setup} gives the experimental protocol. Section~\ref{sec:results} presents results, and Section~\ref{sec:discussion} discusses their clinical meaning, limitations and ethical aspects. Section~\ref{sec:conclusion} concludes.
